\section{Resumen}
Se presenta el desarrollo del proyecto ``Sistema Web de Gestión Documental y Generación de Códigos QR para el PJETAM'' realizado en el Poder Judicial del Estado de Tamaulipas ubicado en Boulevard Praxedis Balboa número 2207, entre las calles López Velarde y Díaz Mirón, Colonia Miguel Hidalgo, CP 87090, en Ciudad Victoria, Tamaulipas.
El objetivo principal del proyecto consistió en desarrollar una aplicación web que permitiera al personal, convertir, almacenar y gestionar documentos e imágenes, generando códigos QR vinculados a cada archivo para facilitar su referenciación. La aplicación se integró con los sistemas de gestión de escritorio existentes mediante autenticación por parámetro cifrado en la URL.

El sistema se creó utilizando la arquitectura MVC sobre ASP.NET Core 10 con el lenguaje de programación C\#, SQL Server para la base de datos gestionado mediante Entity Framework Core, AWS S3 para el almacenamiento de archivos en la nube, para los servicios de conversión a PDF se utilizó LibreOffice en modo headless y para el de generación de códigos QR la librería QRCoder. Para el desarrollo de las interfaces de usuario se emplearon vistas Razor (.cshtml) combinadas con Bootstrap 5. 
Los conocimientos adquiridos durante la educación superior en bases de datos relacionales, almacenamiento en la nube, programación orientada a objetos y desarrollo web permitieron unir diferentes tecnologías para desarrollar un sistema completo y funcional. 

Al finalizar el proyecto se obtuvo el módulo de Conversión y Carga de archivos y el módulo de Gestión Documental, que centralizó la consulta, visualización, descarga, eliminación y generación de códigos QR. De esta manera, el sistema cumplió con los requerimientos establecidos, una sola aplicación con todos los procesos de gestión documental para eliminar las herramientas externas y procesos lentos.

%  que anteriormente se hacían de forma manual usando herramientas externas, y mejorando la organización documental de la institución de manera segura y fácil de utilizar.

% que busca resolver este proyecto es la mejora de la gestión documental del Poder Judicial del Estado de Tamaulipas, busca mejorar la digitalización de documentos, eliminar el uso de herramientas externas y facilitar los procesos manuales actuales. 

% El sistema se desarrollo utilizando el patrón de arquitectura MVC, implementado sobre el framework ASP.NET Core 10. El sistema utiliza principalmente el lenguaje C\# para el backend, esto para aprovechar las características disponibles en .NET 10. Para la conversión de documentos se utiliza LibreOffice en modo headless.

% Como motor de base de datos relacional el sistema utiliza SQL Server 2022, para crear la BD se utilizó Entity Framework Core con el enfoque Code First. Para almacenar los archivos el sistema utiliza el servicio AWS S3 garantizando la seguridad de los datos, ya que el bucket utilizado es privado y el acceso se realiza a través del servidor para proteger las credenciales.

% Las mejoras obtenidas tras implementar la solución tecnológica fueron la mejora en la organización y gestión de los documentos, obteniendo la centralización de los documentos, y eliminando el uso de herramientas externas. Así como la trazabilidad de acciones realizadas por cada área mediante los registros en la base de datos.

% El aporte principal del proyecto es la generación automática de códigos QR vinculados a cada archivo y la centralización de acciones disponibles, ya que el sistema almacena, convierte si es necesario, y permite visualizar, descargar, eliminar los archivos. 

% Palabras clave
% Gestión documental, código QR, conversión PDF, AWS S3, ASP.NET Core, MVC,